\documentclass[11pt]{article}
\usepackage[margin=1in]{geometry}
\usepackage{times}
\usepackage{hyperref}
\hypersetup{colorlinks=true, linkcolor=black, urlcolor=blue, citecolor=black}
\title{Ionized Air and Borrowed Momentum: The Biefeld-Brown Effect as a Gravitational Phenomenon}
\author{Flyxion \\ \small Independent Researcher}
\date{}
\begin{document}
\maketitle

\begin{abstract}
The Biefeld-Brown effect, in which an asymmetric high-voltage capacitor experiences a net thrust toward its smaller, sharper electrode, has been claimed since the early twentieth century as evidence of electrical control over gravitation. This paper reviews the effect against a field-theoretic and entropic account of gravity in which spacetime curvature tracks the redistribution of an ordered configuration rather than the action of a discrete force, and argues that the observed thrust is accounted for by corona discharge and ionic wind, a conventional electrohydrodynamic phenomenon with no gravitational content. The paper also treats the one serious complication in the record, an Army Research Laboratory estimate that early ion wind models fell short of the measured force, and shows how later work removed that shortfall. The persistence of the gravitational interpretation is treated as a case study in mistaking a correlation between an apparatus's orientation and its motion for a causal channel that does not exist. The paper does not claim to have measured the force in any new apparatus; it assembles the published evidence and states where that evidence is thin.
\end{abstract}

\section{Introduction}

An asymmetric capacitor charged to several tens of kilovolts will, if suspended freely, drift toward its smaller electrode. The observation is credited to Thomas Townsend Brown and his teacher Paul Alfred Biefeld in the 1920s, and Brown secured a British patent on a method of producing force or motion from it in 1928 \cite{Brown1928}. The observation itself is not in dispute. What has been in dispute for a century is the mechanism, with one camp attributing the motion to a direct electrical coupling to the gravitational field and a wider consensus attributing it to the movement of air ionized at the sharp electrode and accelerated toward the blunt one, carrying the apparatus along by momentum exchange with the surrounding atmosphere.

The gravitational reading survives less on evidence than on the intuitive appeal of a device that appears to move without a visible mechanical connection to anything pushing it. The appeal predates a great deal of twentieth century engineering and has never been a reliable guide to what is actually happening. It was reinforced by Brown's later patents, which describe the force as occurring in vacuum and which therefore carry the claim that the thrust cannot be a reaction against air \cite{Brown1965}. Whether that claim holds is the empirical question on which the whole matter depends, and the record on it is more informative than the usual summaries suggest.

\section{What Gravity Is, on the Account Assumed Here}

Rather than treat gravity as a discrete attractive force transmitted between masses, the account adopted in this paper treats gravitational curvature as tracking the redistribution of an ordered configuration, the way entropy gradients in a coupled scalar, vector, and entropy field organize the geometry a test body moves through. On this view there is no local switch that couples to gravitation the way a solenoid couples to a magnetic field, because gravitation is not a locally sourced force in the relevant sense; it is a large-scale bookkeeping of how ordered configurations are arranged and how that arrangement constrains motion. An electrical apparatus, however cleverly built, has no channel into that bookkeeping unless it is moving enough mass-energy to matter at cosmological or at least planetary scale, which a laboratory capacitor emphatically is not.

This matters because it removes, from the outset, any plausible mechanism by which a kilovolt potential across two electrodes could locally alter the surrounding curvature. Whatever an asymmetric capacitor is doing, it is not doing it by touching the gravitational field, because there is no local handle on that field to touch. The argument does not need this framework to go through. In ordinary general relativity the energy stored in a capacitor of laboratory size gravitates in proportion to its mass-energy, a quantity many orders of magnitude too small to produce a thrust of the size observed, so a gravitational reading of the effect would need a coupling that general relativity does not contain.

\section{What Is Actually Happening}

The conventional account is unglamorous and, in its main lines, complete. High voltage at a sharp electrode ionizes the surrounding air through corona discharge. The resulting ions are accelerated by the electric field toward the oppositely charged, typically blunter electrode, colliding with neutral air molecules along the way and dragging them into a bulk flow, the ionic wind. By Newton's third law, the apparatus experiences a reaction force opposite to the direction of the ejected air, which is toward the sharp electrode, the direction in which the device is observed to move. This was already the proposed explanation in the 1960s, when Christenson and Moller built electrode configurations of the Biefeld-Brown type, measured the thrust in ambient air, and found good agreement with a theory of ion-neutral momentum transfer \cite{CM1967}.

The distinguishing prediction is dependence on the gas. A reaction against ionized air should change with pressure and humidity and should vanish when the gas is removed. A coupling to gravity, on the account assumed here, gives no reason for the presence of air to matter. Tajmar reviewed the claims and theory and ran a dedicated experiment to look for any thrust not attributable to corona discharge, and concluded that any such contribution would have to lie at least five orders of magnitude below the corona wind effect, which he read as showing that the effect is a misinterpretation of corona wind phenomena and not a link between gravitation and electromagnetism \cite{T2004}. Earlier vacuum testing pointed the same way: Talley's tests of Brown's electrode configurations at high vacuum reported that no detectable propulsive force was electrostatically induced by a static potential across the electrodes \cite{Talley1991}.

\section{The Army Research Laboratory Complication}

An honest account has to deal with the report that did the most to keep the question open. Bahder and Fazi of the Army Research Laboratory built several asymmetric capacitors, confirmed that the force is real, and then estimated the size of the force under two models of charge conduction. They found that a ballistic ionic wind model gave a force at least three orders of magnitude too small to explain the measured value, while an ion drift scaling theory was offered as a more promising route, and they concluded that the physical basis of the effect was not at that time understood \cite{BF2003}. That sentence has been quoted in the antigravity literature as if it were a finding of anomalous physics. It is better read for what it says: a particular simple estimate of the momentum carried by the wind was inadequate, and the calculation needed to be done better.

Subsequent work did that. Canning, Melcher and Winet tested several asymmetrical capacitor thrusters for NASA and found that the thrust depended on voltage, polarity and grounding configuration in a way consistent with ion motion, and that devices without an accompanying ion wind produced a null thrust \cite{CMW2004}. Ianconescu, Sohar and Mudrik then computed the lifting force from the unicharge ion wind and presented experimental results that confirmed the theoretical values \cite{ISM2011}. The strongest demonstration is practical. Xu and colleagues flew a fixed-wing aircraft across a gymnasium with no moving parts, using ionic wind thrust, an application that would be impossible if the thrust came from anything other than momentum given to the air \cite{Xu2018}.

\section{Objections}

Proponents sometimes note that thrust in near-vacuum, though reduced, does not always vanish entirely, and treat the residual as evidence of a genuine field effect beneath the ionic wind. The more parsimonious explanation is outgassing and residual ionization from imperfect vacuum, plus leakage currents through whatever the apparatus is mounted on, both of which are difficult to eliminate completely and neither of which requires appeal to gravitation. Brown's own 1965 patent describes the force as occurring even in vacuum, and an assertion in a patent specification is not a controlled measurement; the measurements that exist at comparable pressures do not show it \cite{Brown1965,Talley1991}.

A second objection holds that even if ionic wind explains most of the effect, some fraction remains unaccounted for and that fraction is where the gravitational coupling lives. This is an argument from the residual of an imperfect experiment rather than from a positive detection. A residual of this kind has to show structure to count as a signal: it should grow rather than shrink as measurement improves, it should vary with the parameters a gravitational coupling would depend on, and it should be seen independently by groups with different apparatus. The published record does not show it. The residual claimed in vacuum has shrunk toward the noise floor as techniques improved, which is the pattern expected of experimental error rather than of a real, if small, effect.

A third objection points to the original Army report and asks why a laboratory with no stake in the matter wrote that the effect was not understood. The answer is that the sentence described a gap in a calculation, which was then closed. The report also recommended that the effect be tested in vacuum, which is the experiment on which the gravitational reading has consistently failed \cite{BF2003}.

\section{Why the Gravitational Reading Persists}

Three features of the effect favor its misreading. The device moves silently, with no jet, propeller or visible exhaust, and the working fluid is the invisible air around it. The thrust is directional and reproducible, which is the signature of a real force and so seems to demand a real, novel cause. And the effect has a founding myth: a capacitor builder's observation of motion in the 1920s, followed by decades of patents describing it in the language of gravity control, gives the claim a lineage that the ordinary electrohydrodynamic account lacks. None of these features is evidence. Each is a reason for the claim to circulate.

\section{Conclusion}

The Biefeld-Brown effect is real, well characterized, and electrohydrodynamic. Its persistence as an antigravity claim owes more to the emotional appeal of an apparatus that moves without visible mechanical linkage than to any surviving evidentiary gap. The one serious complication in the record, the Army estimate that the simple wind model undershot the force, has been resolved by better modeling and by direct tests. On an account of gravity as the bookkeeping of large-scale ordered configuration rather than a locally sourced force, there was never a mechanism available for a benchtop capacitor to exploit in the first place, which makes the ionic wind explanation not merely more parsimonious but the only one with a channel to act through. This paper offers no new measurement and no positive theory of gravitational coupling; it states only that the evidence for such coupling is absent and that the evidence for the conventional account is substantial.

\begin{thebibliography}{99}

\bibitem{Brown1928} T. T. Brown, A Method of Producing Force or Motion, British Patent 300,311, 1928.

\bibitem{Brown1965} T. T. Brown, Electrokinetic Apparatus, U.S. Patent 3,187,206, issued June 1, 1965.

\bibitem{CM1967} E. A. Christenson and P. S. Moller, Ion-Neutral Propulsion in Atmospheric Media, \emph{AIAA Journal} 5 (1967) 1768--1773.

\bibitem{Talley1991} R. L. Talley, Veritay Technology Inc., test report on electrostatically induced propulsion at high vacuum, prepared for the U.S. Air Force, 1991, as discussed in the literature on asymmetrical capacitor thrusters.

\bibitem{BF2003} T. B. Bahder and C. Fazi, Force on an Asymmetric Capacitor, Army Research Laboratory Technical Report ARL-TR-3005, June 2003, arXiv:physics/0211001.

\bibitem{CMW2004} F. X. Canning, C. Melcher and E. Winet, Asymmetrical Capacitors for Propulsion, NASA Contractor Report NASA/CR-2004-213312, October 2004.

\bibitem{T2004} M. Tajmar, Biefeld-Brown Effect: Misinterpretation of Corona Wind Phenomena, \emph{AIAA Journal} 42 (2004) 315--318.

\bibitem{ISM2011} R. Ianconescu, D. Sohar and M. Mudrik, An Analysis of the Brown-Biefeld Effect, \emph{Journal of Electrostatics} 69 (2011) 512--521, arXiv:1011.1393.

\bibitem{Xu2018} H. Xu, Y. He, K. L. Strobel, C. K. Gilmore, S. P. Kelley, C. C. Hennick, T. Sebastian, M. R. Woolston, D. J. Perreault and S. R. H. Barrett, Flight of an Aeroplane with Solid-State Propulsion, \emph{Nature} 563 (2018) 532--535.

\end{thebibliography}

\end{document}
